% Artigo técnico — UCDB / Redes
% pdfLaTeX + babel português
\documentclass[12pt,a4paper]{article}

\usepackage[utf8]{inputenc}
\usepackage[T1]{fontenc}
\usepackage[brazil]{babel}
\usepackage{graphicx}
\usepackage{hyperref}
\usepackage{booktabs}
\usepackage{amsmath}

\title{Otimização de APIs REST e cache HTTP para exposição de datasets massivos}
\author{Seu Nome \\ Universidade Católica Dom Bosco (UCDB) — Redes}
\date{2026}

\begin{document}
\maketitle

\begin{abstract}
Este trabalho aborda o problema de disponibilizar, via API REST, conjuntos de dados da ordem de milhões de registros---cenário análogo à publicação de bases estatísticas (por exemplo, dados do IBGE)---sem comprometer estabilidade nem latência do servidor. São discutidas paginação por \emph{offset} versus paginação por \emph{cursor}, compressão algorítmica (Gzip e Brotli), políticas de cache em cabeçalhos (\texttt{Cache-Control}) e redução da verbosidade de payloads JSON. Descreve-se uma implementação em FastAPI com SQLite, documentação OpenAPI, middleware de compressão Gzip e validação condicional com \texttt{ETag} e respostas \texttt{304 Not Modified}, de forma que clientes não retransmitam páginas inalteradas após mudanças globais versionadas do dataset.
\end{abstract}

\section{Introdução}

Serviços que expõem grandes volumes tabulares pela Web enfrentam limites físicos de memória, CPU e largura de banda. Enviar ``tudo de uma vez'' viola boas práticas REST e tende a causar timeouts, \emph{spikes} de uso de RAM e custo elevado de serialização JSON. A solução usual combina \textbf{paginação}, \textbf{filtros eficientes}, \textbf{compressão na camada HTTP} e \textbf{cache com revalidação} (\emph{conditional requests}).

\section{Fundamentação teórica}

\subsection{Paginação: \emph{offset} versus \emph{cursor}}

A abordagem clássica \texttt{LIMIT $n$ OFFSET $k$}, em muitos SGBDs, exige percorrer ou descartar internamente as primeiras $k$ linhas para entregar a página seguinte, com custo que cresce com $k$. Em bases com milhões de linhas, páginas profundas degradam o tempo de resposta e aumentam carga no disco.

A paginação por \textbf{cursor} utiliza um marcador opaco (tipicamente a chave da última linha retornada) e uma cláusula do tipo ``\texttt{WHERE id > :cursor ORDER BY id LIMIT :n}'', permitindo uso de índice e tempo aproximadamente constante por página. O trade-off é menor ergonomia para ``pular'' arbitrariamente para a página 10\,000 sem percorrer sequência de cursores; em APIs de dados públicos, cursor é frequentemente preferível para leitura sequencial e \emph{streaming} de consumo.

\subsection{Compressão: Gzip e Brotli}

JSON textual apresenta alta redundância; algoritmos como \textbf{Gzip} e \textbf{Brotli} reduzem o número de octetos transmitidos. Há custo de CPU no servidor e no cliente. Em redes lentas ou payloads grandes, a redução de bytes costuma compensar o custo de compressão; em latências muito baixas e corpos pequenos, compressão pode não valer a pena. Em ambientes de produção, Brotli é comum atrás de proxies (nginx, CDN); frameworks minimalistas costumam oferecer Gzip nativamente.

\subsection{Cabeçalhos de cache: \texttt{Cache-Control}}

O cabeçalho \texttt{Cache-Control} orienta caches intermediários (CDN, proxies) e o próprio navegador. Valores como \texttt{private} impedem armazenamento compartilhado indevido de dados sensíveis; \texttt{max-age} define frescor sem revalidação. Para dados que mudam ocasionalmente, combina-se \emph{freshness} curta com \textbf{revalidação condicional} (\texttt{ETag} / \texttt{If-None-Match}), permitindo respostas \texttt{304 Not Modified} sem reenvio do corpo.

\subsection{Verbosidade do JSON}

Nomes de campos longos e estruturas aninhadas sem necessidade aumentam bytes e tempo de parse. Uma modalidade é oferecer representações \textbf{mínimas} (chaves curtas) para consumo programático e representações \textbf{verbosas} para legibilidade humana ou depuração, selecionáveis por parâmetro de consulta, mantendo o mesmo contrato lógico na camada de domínio.

\section{Metodologia}

A solução usa \textbf{FastAPI} e \textbf{SQLite}. Para testes de volume elevado, pode-se carregar milhões de registros sintéticos (identificador, código estilo IBGE, nome, UF e metadados); para dados oficiais, a API dos municípios do IBGE alimenta a base com cerca de 5,5 mil registros. A tabela possui índice primário em \texttt{id} e índice composto \texttt{(uf, id)} para acelerar filtros por unidade federativa combinados à paginação.

O endpoint principal \texttt{GET /items} aceita \texttt{limit}, \texttt{cursor}, \texttt{uf}, \texttt{q} (prefixo textual para uso de índice) e \texttt{fields=minimal|full}. A documentação \textbf{OpenAPI} é gerada automaticamente (\texttt{/docs}, \texttt{/openapi.json}).

Compressão: a API negocia \textbf{Brotli} (\texttt{br}) ou \textbf{Gzip} (\texttt{gzip}) via \texttt{Accept-Encoding}, com limiar mínimo de bytes. Para comparação acadêmica com \textbf{OFFSET}, um endpoint opcional \texttt{/items/debug/offset-page} fica desabilitado por padrão e exige variável de ambiente \texttt{ENABLE\_BENCHMARKS=1}.

\subsection{ETag e validação em larga escala}

Calcular checksum sobre milhões de linhas a cada requisição é inviável. Adota-se uma linha em \texttt{dataset\_meta} com inteiro \texttt{version}, bumpado quando o dataset é recarregado ou atualizado globalmente. O \texttt{ETag} de uma página listada é o hash SHA-256 de uma tupla canônica $(\texttt{version}, \text{parâmetros da consulta})$. Assim, a verificação de \texttt{If-None-Match} custa uma leitura da meta e hash, sem varrer a base. Quando o cliente reenvia o mesmo \texttt{ETag}, a API responde \texttt{304} \textbf{sem executar} a consulta paginada nem serializar JSON.

Para um item individual, o \texttt{ETag} incorpora também o carimbo \texttt{updated\_at}, garantindo invalidação se o registro mudar.

\section{Resultados esperados e medição}

Com paginação por cursor e índices adequados, espera-se latência estável por página mesmo com milhões de linhas na tabela. Com \texttt{OFFSET} elevado, espera-se crescimento marcante do tempo de consulta (medido na aplicação pelo cabeçalho \texttt{X-Query-Time-Ms} e repetível conforme roteiro em \texttt{docs/METRICAS.md}).

A compressão (Gzip/Brotli) deve reduzir o tamanho transferido em listas com \texttt{fields=full}. A segunda requisição idêntica com \texttt{If-None-Match} deve produzir \texttt{304} e corpo vazio, economizando banda proporcional ao tamanho da página.

\begin{table}[h]
\centering
\begin{tabular}{@{}ll@{}}
\toprule
Experimento & Indicador sugerido \\
\midrule
Cursor vs OFFSET alto & Tempo de consulta (ms) \\
Gzip/Brotli on/off & Bytes baixados (\texttt{curl} / ferramentas) \\
304 & Presença de corpo vazio e código HTTP \\
\bottomrule
\end{tabular}
\caption{Experimentos para gráficos do artigo (preencher com medições reais).}
\end{table}

\section{Conclusão}

A combinação de paginação por cursor, filtros indexáveis, compressão HTTP, políticas de \texttt{Cache-Control} e validação condicional com \texttt{ETag} permite servir datasets massivos de forma sustentável. A implementação mostra que \texttt{304 Not Modified} em escala exige \textbf{metadados de versão} e \textbf{canonicalização de parâmetros}, não recomputação sobre todo o conjunto.

\begin{thebibliography}{9}
\bibitem{fielding}
Fielding, R. T. \emph{Architectural Styles and the Design of Network-based Software Architectures}. Tese de doutorado, UC Irvine, 2000.

\bibitem{rfc9110}
Fielding, R., Nottingham, M. \emph{HTTP Semantics}. RFC 9110, IETF, 2022. Disponível em: \url{https://www.rfc-editor.org/rfc/rfc9110}

\bibitem{fastapi}
Ramírez, S. \emph{FastAPI}. Documentação oficial. \url{https://fastapi.tiangolo.com/}
\end{thebibliography}

\end{document}
